\honest{Bystander Apathy}{Eliezer Yudkowsky}{2009}

Larger groups are less likely to act in an emergency, and not only one by one: taken together. Put a subject alone in a room with smoke coming in, and 75\% report it. Put three real subjects together, and in only 38\% of the groups does anyone report it. With two confederates who ignore the smoke, 10\% report it. \nb{The figures match Latané and Darley's article. The samples are small: 24 subjects alone, ten with confederates, eight groups of three. A 2020 study of real public conflicts filmed by CCTV found that someone intervened in nine of ten cases, and that more bystanders went with a greater chance that someone would.} \nb{Latané and Darley's article opens with the 1964 murder of Kitty Genovese and the thirty-eight neighbors said to have watched in silence. The post does not mention the case. Later research found no evidence for 38 witnesses or for their inaction, and in 2016 the New York Times said its 1964 report had ``grossly exaggerated the number of witnesses and what they had perceived.''}

The standard explanations are diffusion of responsibility, in which everyone hopes someone else will bear the cost of acting first and the crowd offers an excuse, and pluralistic ignorance: everyone tries to look calm while looking for cues, and sees that everyone else looks calm. I quote Cialdini on this: an emergency is often not obviously one (is the man in the alley having a heart attack, or sleeping off a drink?), so we look to others for clues, but placidly, with camouflaged glances, and see everyone unruffled. Cialdini's advice: if you need help, point at one person and ask that person, since the whole group may be less likely to help than one individual.

Could evolution have produced an arms race to be the last to help? In a band of relatives, the first to act pays the cost while the others share the genetic benefit. I think not: once the whole group fails to act, a gene for helping at once should invade, and the experiments show people failing to help, not waiting. \nb{The arms race was the author's own guess, in ``Evolving to Extinction'' (2007). Here the author rejects it.} More likely, I think, the problem is that the bystanders are strangers: helping goes up when subjects know the victim, and, ``If I recall correctly,'' when they know each other. \nb{Both findings are in the article the post cites. Knowing the victim meant a chat of less than a minute in the hall.}

Nervousness may play a part too: if Robin Hanson is right about ``choking,'' being first to act might be a dangerous bid for status. I cannot recall seeing shyness discussed in analyses of the effect, ``but that's probably just my poor memory.'' \nb{The cited article discusses bystanders as ``an audience'' to each other, trying ``to avoid possible ridicule and embarrassment.''}

Cynically, people may care mostly about not being blamed. That may contribute, but it is not the whole story. Subjects ignored smoke that could have threatened them, and telling people about the bystander effect reduces it. \nb{The post gives no source. A 1978 study by Beaman and colleagues supports it.} This is one of the few biases, as I recall, that teaching seems to cure, perhaps because the correction is obvious and hard to overdo, unlike adjusting your calibration. People sometimes hold themselves responsible, once they see they are the only ones who know enough to act. \nb{Latané and Darley agree: the people who did not help were in conflict, not indifferent. Their title puts ``Apathy'' in quotation marks.}

I end by wondering what happens in a crowd where everyone has been told about the effect.
